\section*{Hoofdstuk \ref{ch:g10:freshwater-resources} --- Zoetwatervoorraden}

\begin{solution}{exo:g10:freshwater-resources:1}
Een \omterm{def:g10:freshwater-resources:stock}{voorraad} is de hoeveelheid water die op een bepaald moment in een reservoir
zit, in \unit{km^3} (of kilogram). Een \omterm{def:g10:freshwater-resources:stock}{flux} is de hoeveelheid die per
tijdseenheid van het ene reservoir naar het andere gaat, in \unit{km^3} per jaar.
De \omterm{def:g10:freshwater-resources:residence-time}{verblijftijd} is de gemiddelde tijd die water in een reservoir doorbrengt, in
dagen of jaren; in stationaire toestand is ze de \omterm{def:g10:freshwater-resources:stock}{voorraad} gedeeld door de \omterm{def:g10:freshwater-resources:stock}{flux}.
\end{solution}

\begin{solution}{exo:g10:freshwater-resources:2}
Ongeveer \qty{2.5}{\percent} van al het water is zoet. Van het zoete water is
\qty{68.7}{\percent} bevroren in ijskappen, gletsjers en eeuwige sneeuw, en zit
\qty{30.1}{\percent} onder de grond. De rivieren bevatten het minst,
\num{2120}~\unit{km^3}.
\end{solution}

\begin{solution}{exo:g10:freshwater-resources:3}
\omterm{def:g10:freshwater-resources:renewable}{Hernieuwbaar zoet water} is het zoete water dat elk jaar door neerslag wordt
vernieuwd en als afstroming en grondwater door een gebied stroomt: een volume per
jaar, een \omterm{def:g10:freshwater-resources:stock}{flux}. \omterm{def:g10:freshwater-resources:fossil-groundwater}{Fossiel grondwater} is water dat lang geleden in een watervoerende
laag is gekomen en dat de huidige aanvulling niet vernieuwt: het is een \omterm{def:g10:freshwater-resources:stock}{voorraad}
die alleen kan worden uitgeput, een stock.
\end{solution}

\begin{solution}{exo:g10:freshwater-resources:4}
$91\,000/2120 \approx 43$ keer meer.
\end{solution}

\begin{solution}{exo:g10:freshwater-resources:5}
De atmosfeer vernieuwt haar water het snelst (ongeveer \num{9.5} dagen), de
oceaan het traagst van de berekende reservoirs (ongeveer \num{3200} jaar); diep
grondwater is nog ouder. $3200 \times 365.25/9.5 \approx \num{120000}$
keer langer.
\end{solution}

\begin{solution}{exo:g10:freshwater-resources:6}
Erin: $413\,000 + 73\,000 = \num{486000}$~\unit{km^3} per jaar, dus
$T = 12\,700/486\,000 \approx \num{0.0261}$ jaar $\approx \num{9.5}$ dagen.
Eruit: $373\,000 + 113\,000 = \num{486000}$~\unit{km^3}, evenveel: de lucht is in
stationaire toestand, en de \omterm{def:g10:freshwater-resources:residence-time}{verblijftijd} is dezelfde.
\end{solution}

\begin{solution}{exo:g10:freshwater-resources:7}
$1\,335\,040/413 \approx \num{3230}$ jaar. In stationaire toestand ontvangt de
oceaan evenveel als hij verliest: $373 + 40 = 413$ duizend~\unit{km^3} per jaar,
dus de \omterm{def:g10:freshwater-resources:residence-time}{verblijftijd} berekend met de inkomsten is dezelfde.
\end{solution}

\begin{solution}{exo:g10:freshwater-resources:8}
Het land ontvangt \num{113000}~\unit{km^3} per jaar; \num{73000} keren terug naar
de lucht door verdamping en transpiratie en \num{40000} stromen af naar zee:
$73\,000 + 40\,000 = \num{113000}$. Het deel dat naar de lucht terugkeert, is
$73/113 \approx \qty{65}{\percent}$.
\end{solution}

\begin{solution}{exo:g10:freshwater-resources:9}
$273.2 \times 10^6 \times 1233.5 \approx \qty{3.37e11}{m^3}$, dat is
\qty{337}{km^3}: bijna evenveel als al het water dat het Aralmeer in 1987 nog
bevatte.
\end{solution}

\begin{solution}{exo:g10:freshwater-resources:10}
$15.8 \times 0.3048 \approx \qty{4.82}{m}$. Over de 65 jaar van 1950 tot
2015 is dat $482/65 \approx \qty{7.4}{cm}$ per jaar.
\end{solution}

\begin{solution}{exo:g10:freshwater-resources:11}
$2/0.01 = 200$ jaar. De \omterm{def:g10:freshwater-resources:stock}{voorraad} daalt nu met $0.05 - 0.01 =
\qty{0.04}{km^3}$ per jaar, dus hij gaat $2/0.04 = 50$ jaar mee.
\end{solution}

\begin{solution}{exo:g10:freshwater-resources:12}
1960: $1093/71.1 \approx 15$ jaar; 2010: $98.1/11.4 \approx 8.6$ jaar.
De verdamping neemt elk jaar ongeveer dezelfde waterdikte weg van elke vierkante
kilometer van het meer; toen het meer kromp, kromp zijn oppervlakte, en daarmee
ook het totale verdampte volume.
\end{solution}

\begin{solution}{exo:g10:freshwater-resources:13}
(1) Het nettoverlies is $1.2 - 0.2 = \qty{1.0}{m}$ per jaar over
$100 \times 10^6$~\unit{m^2}: $10^8$~\unit{m^3}, er moet \qty{0.1}{km^3} per jaar
instromen. (2) De toevoer zakt tot \qty{0.05}{km^3}: het meer verliest
$\qty{0.05}{km^3} = 5 \times 10^7$~\unit{m^3}, een laag van
$5 \times 10^7/10^8 = \qty{0.5}{m}$. (3) Naarmate het peil daalt, komen de
oevers boven water en krimpt de oppervlakte van het meer: minder oppervlak,
minder verdamping, tot verdamping en toevoer weer in evenwicht zijn.
\end{solution}

\begin{solution}{exo:g10:freshwater-resources:14}
De watervoerende laag bevat $100 \times 0.20 = \qty{20}{m}$ water, dat is
\qty{20000}{mm}: met \qty{1}{mm} per jaar \num{20000} jaar. Het diepe water van
de watervoerende laag onder de Sahara is \num{200000} tot een miljoen jaar oud,
tien tot vijftig keer meer: het is aangevuld tijdens nattere klimaten in het
verleden.
\end{solution}

\begin{solution}{exo:g10:freshwater-resources:15}
$337/453\,000 \approx \qty{0.00074}{km}$, ongeveer \qty{0.74}{m} water. Het
waterpeil daalde \qty{4.8}{m}, ongeveer zes en een half keer meer. Het water vult
alleen de poriën van het gesteente (en loopt er niet helemaal uit): wanneer de
grondwaterspiegel een meter daalt, is het water dat vrijkomt een laag die veel
dunner is dan een meter, hier ongeveer $0.74/4.8 \approx \num{0.15}$~m.
\end{solution}

\begin{solution}{pb:g10:freshwater-resources:1}
\textbf{1.} $1093/68\,000 \approx \qty{0.016}{km}$, ongeveer \qty{16}{m}.

\textbf{2.} $1093 \times 10^9 \approx \qty{1.09e12}{m^3}$.

\textbf{3.} \qty{10}{g/L} is \qty{10}{kg/m^3}: $1.093 \times 10^{12} \times 10
\approx \qty{1.1e13}{kg}$, ongeveer \num{11} miljard ton.

\textbf{4.} $1093/71.1 \approx 15$ jaar.

\textbf{5.} Er stroomt geen rivier uit: het water kan er alleen als damp uit
(of door in de grond weg te sijpelen), terwijl het zout achterblijft.

\textbf{6.} $71.1/68\,000 \approx \qty{0.00105}{km}$, ongeveer \qty{1.05}{m}.

\textbf{7.} $9.4/68\,000 \approx \qty{0.00014}{km}$, ongeveer \qty{0.14}{m}.

\textbf{8.} $71.1 - 9.4 = \qty{61.7}{km^3}$ per jaar.

\textbf{9.} $79 + 37 = \qty{116}{km^3}$: $61.7/116 \approx \qty{53}{\percent}$,
ongeveer de helft.

\textbf{10.} Inkomsten $1.0 + 6.2 = \qty{7.2}{km^3}$; uitgaven \qty{36.8}{km^3};
nettoverandering $7.2 - 36.8 = \qty{-29.6}{km^3}$.

\textbf{11.} $(502.7 - 327.2)/5 = \qty{35.1}{km^3}$ per jaar: dezelfde orde van
grootte als de \qty{29.6}{km^3} van de balans van 1987 (afzonderlijke jaren
verschillen: de verdamping was \qty{47.9}{km^3} in 1984).

\textbf{12.} $36.8/0.00105 \approx \qty{35000}{km^2}$.

\textbf{13.} $35\,000/68\,000 \approx \qty{52}{\percent}$, ongeveer de helft.

\textbf{14.} Het meer was gekrompen: de helft van het oppervlak verdampt ongeveer
de helft van het water.

\textbf{15.} $345.6/1093 \approx \qty{32}{\percent}$, ongeveer een derde.

\textbf{16.} $10 \times 1093/345.6 \approx \qty{31.6}{g/L}$.

\textbf{17.} Het gemeten zoutgehalte is lager: een deel van het zout had het water
verlaten, was gekristalliseerd op de drooggevallen bodem en door de wind als zout
stof weggevoerd.

\textbf{18.} Ongeveer \qty{1.05}{m} per jaar, de verdampte dikte.

\textbf{19.} $7.2/35\,000 \approx \qty{0.0002}{km}$, ongeveer \qty{0.20}{m}.

\textbf{20.} $1.05 - 0.20 \approx 0.85$ (\num{0.84} met de niet-afgeronde waarden):
\textbf{in 1987 daalde het peil van het Aralmeer met ongeveer \qty{0.8}{m}, bijna
een meter in één jaar}.
\end{solution}
